% Critical Field Analysis
% Source: docs/critical_field.md (migrated to LaTeX)
% Uses macros from preamble.tex

\section{Critical Field Analysis}
\label{sec:critical-field}

\textbf{Purpose}: Determine the minimum magnetic field strength required for full
graviton-photon conversion in a given physical setup, and compare conversion
efficiency between theories.

\subsection{Physics Motivation}
\label{sec:cf-motivation}

The Gertsenshtein effect converts gravitons to photons (and vice versa) in a
background magnetic field. The key experimental question is: \textbf{how strong
must the magnetic field be to achieve full conversion?}

For the standard Einstein--Maxwell (E-M) theory, the conversion probability
after a wavepacket transits a localized B-field region is given by the
Boccaletti formula:
\begin{equation}
  P_\text{final} = \sin^2\!\left(\frac{\kappa}{2} \int B(z)\, dz\right)
  \label{eq:cf-boccaletti}
\end{equation}

Full conversion ($P_\text{final} = 1$) occurs at the first peak of the $\sin^2$
oscillation, when:
\begin{equation}
  \frac{\kappa}{2} \int B(z)\, dz = \frac{\pi}{2}
  \label{eq:cf-full-conversion}
\end{equation}

For a Gaussian profile $B(z) = B_\text{peak} \exp\!\bigl(-z^2 / 2R^2\bigr)$, this gives:
\begin{equation}
  B_\text{peak}^\text{min} = \frac{\pi}{\kappa\, R\, \sqrt{\pi/2}}
  \label{eq:Bpeak-min}
\end{equation}

For theories with additional couplings (torsion, Chern--Simons, axion-photon),
the conversion formula differs. The critical field analysis finds $B_\text{min}$
numerically from sweep data, enabling direct comparison between theories.

\subsection{Key Concepts}
\label{sec:cf-concepts}

\subsubsection{Why $P_\text{final}$, not $P_\text{max}$}
\label{sec:pfinal-vs-pmax}

During transit through a localized B-field, energy oscillates (Rabi
oscillation) between graviton and photon. A momentary $P = 1$ inside the
field region is \textbf{not} full conversion --- the energy may oscillate back
before exit.

The physically meaningful quantity is $\bm{P_\text{final}}$: the permanent net
conversion after the wavepacket exits the B-field region. This is what a
downstream detector would measure, and what the Boccaletti formula predicts.

For uniform $B$ in a periodic domain, $P_\text{final} = P(t_\text{end}) = \sin^2(\kappa\, \Bzero\, t_\text{end} / 2)$, which is the conversion at the observation time.

\subsubsection{Amplification Factor}
\label{sec:amplification}

To compare theories, define the amplification factor:
\begin{equation}
  \text{amplification} = \frac{B_\text{min}(\text{E-M})}{B_\text{min}(\text{new theory})}
  \label{eq:amplification}
\end{equation}
at the same threshold ($P_\text{final} \geq 0.99$). If:
\begin{itemize}
  \item amplification $> 1$: new theory enhances conversion (needs less $B$)
  \item amplification $< 1$: new theory weakens conversion
  \item amplification $= 1$: identical to E-M
\end{itemize}

\subsubsection{Threshold}
\label{sec:threshold}

The default threshold is $0.99$ (full conversion within numerical tolerance).
Regions of parameter space where no $B$ field achieves $P_\text{final} \geq 0.99$
produce $B_\text{min} = \text{NaN}$, shown as blank spots on heatmaps --- which is
itself informative.

\paragraph{Regime validity caveat.}
For the Gertsenshtein effect, $P_\text{final} = 0.99$ is \textbf{unphysical}: the
mixing length $L_\text{mix} \approx 2\,\text{Mpc}/B_0\,[\text{G}]$ ensures that no
astrophysical or laboratory configuration achieves $P \sim 1$
(see~\S\ref{subsec:physics-regimes} of \texttt{gertsenshtein.tex}).
Moreover, Hwang \& Noh~\cite{hwangnoh2023graviton} show that the flat-background
approximation becomes borderline when $P \sim 1$.

For Gertsenshtein-type sweeps, use $P_0 \leq 0.01$ to remain in the physically
relevant weak-coupling regime.  The amplification factor
$A = B_\text{min}(\text{EM})/B_\text{min}(\text{torsion})$ is
\textbf{threshold-independent} at leading order, since $P \propto B_0^2$ in the
linearized regime (see eq.~\eqref{eq:amplification} in \texttt{gertsenshtein.tex}).
The same amplification can be obtained from the conversion coefficient
$C = P/B_0^2$ without a threshold at all.

\subsection{Algorithm}
\label{sec:cf-algorithm}

Given sweep results with a field-strength parameter (e.g.\ \texttt{B0}) as one of the
swept variables:

\begin{enumerate}
  \item \textbf{Group} rows by unique outer parameter combinations (all swept params
    except the field-strength param)
  \item \textbf{Sort} each group by the field-strength parameter (ascending)
  \item \textbf{Scan} upward for the first crossing where $\text{metric} \geq \text{threshold}$
  \item \textbf{Interpolate} linearly between the two bracketing grid points for
    sub-grid accuracy (optional, enabled by default)
  \item \textbf{Estimate errors} from three sources (see below)
  \item \textbf{Output} a reduced \texttt{SweepResults} with the field-strength parameter
    collapsed, containing \texttt{B\_min}, \texttt{inv\_B\_min}, and error columns
\end{enumerate}

The output plugs directly into the existing \texttt{tidal plot --type sweep}
infrastructure --- no plot code changes needed. 1D line plots, 2D heatmaps,
and 3+ dimension parallel coordinates all work automatically.

\subsection{Error Estimation}
\label{sec:cf-errors}

Three independent error sources, combined in quadrature:

\subsubsection{Grid spacing error (\texttt{B\_min\_err\_grid})}
\label{sec:err-grid}

Always available. Conservative half-interval:
\begin{equation}
  \varepsilon_\text{grid} = \frac{B_\text{hi} - B_\text{lo}}{2}
  \label{eq:err-grid}
\end{equation}

This is the fundamental resolution limit of the sweep grid and dominates for
coarse sweeps.

\subsubsection{Metric propagation error (\texttt{B\_min\_err\_metric})}
\label{sec:err-metric}

Available when ensemble replicates exist. Propagates replicate standard
deviation through the linear interpolation formula using the Jacobian:
\begin{align}
  B_\text{min} &= B_\text{lo} + \frac{T - m_\text{lo}}{m_\text{hi} - m_\text{lo}} \times \Delta B \label{eq:Bmin-interp} \\[6pt]
  \frac{\partial B_\text{min}}{\partial m_\text{lo}} &= -\Delta B \times \frac{m_\text{hi} - T}{(m_\text{hi} - m_\text{lo})^2} \notag \\
  \frac{\partial B_\text{min}}{\partial m_\text{hi}} &= -\Delta B \times \frac{T - m_\text{lo}}{(m_\text{hi} - m_\text{lo})^2} \notag \\[6pt]
  \varepsilon_\text{metric}^2 &= \left(\frac{\partial B_\text{min}}{\partial m_\text{lo}}\right)^2 \sigma_\text{lo}^2 + \left(\frac{\partial B_\text{min}}{\partial m_\text{hi}}\right)^2 \sigma_\text{hi}^2 \label{eq:err-metric}
\end{align}

For $n_\text{replicates} \geq 5$, a more robust alternative: compute $B_\text{min}$ independently
per replicate, take std of the distribution.

\subsubsection{Interpolation model error (\texttt{B\_min\_err\_interp})}
\label{sec:err-interp}

Available when 3+ points exist near the crossing. Compares linear and
quadratic interpolants:
\begin{equation}
  \varepsilon_\text{interp} = |B_\text{min}^\text{linear} - B_\text{min}^\text{quadratic}|
  \label{eq:err-interp}
\end{equation}

The quadratic interpolant uses three points (bracketing pair + one neighbor)
to fit a parabola through the metric values.

\subsubsection{Combined error}
\label{sec:err-combined}

\begin{align}
  \varepsilon_{B_\text{min}} &= \sqrt{\varepsilon_\text{grid}^2 + \varepsilon_\text{metric}^2 + \varepsilon_\text{interp}^2} \label{eq:err-combined} \\
  \varepsilon_{1/B} &= \frac{\varepsilon_{B_\text{min}}}{B_\text{min}^2} \label{eq:err-inv}
\end{align}

\subsubsection{Quality flags}
\label{sec:quality-flags}

Each output point carries a \texttt{crossing\_quality} flag:

\begin{table}[htbp]
\centering
\caption{Crossing quality flags for critical field output.}
\label{tab:quality-flags}
\begin{tabular}{@{}ll@{}}
\toprule
Flag & Meaning \\
\midrule
\texttt{good}   & Clean crossing with interpolation, small errors \\
\texttt{coarse} & Large grid gap at crossing ($\varepsilon_\text{grid} > 0.1 \times B_\text{min}$) \\
\texttt{edge}   & Threshold exceeded at smallest parameter value \\
\texttt{none}   & No crossing found ($B_\text{min} = \text{NaN}$) \\
\bottomrule
\end{tabular}
\end{table}

\subsection{Analytical Reference Formulas}
\label{sec:cf-reference}

Built-in E-M reference formulas compute $B_\text{min}$ analytically for known
geometries, useful for comparison:

\begin{table}[htbp]
\centering
\caption{Built-in analytical reference formulas for $B_\text{min}$.}
\label{tab:reference-formulas}
\begin{tabular}{@{}llll@{}}
\toprule
Formula & Geometry & $P_\text{EM}$ formula & Required params \\
\midrule
\texttt{boccaletti} & Gaussian $B(z)$ & $\sin^2\!\bigl(\kappa/2 \cdot B \cdot R \cdot \sqrt{\pi/2}\bigr)$ & $\kappa$, $R$ \\
\texttt{uniform}    & Uniform $B$, periodic & $\sin^2(\kappa \cdot B \cdot t_\text{end} / 2)$                   & $\kappa$, $t_\text{end}$ \\
\bottomrule
\end{tabular}
\end{table}

For arbitrary geometries (dipolar, 2D/3D), use \texttt{--threshold} with a value
computed from a reference simulation or literature.

\subsection{Usage}
\label{sec:cf-usage}

\subsubsection{Single theory characterization}
\label{sec:single-theory}

\begin{lstlisting}[style=bash]
# 1. Run a parameter sweep including the field-strength parameter
uv run tidal sweep examples/data/gertsenshtein_localized.json \
  --sweep "Bpeak=0.01:0.5:30" --sweep "R=1:20:10" \
  --measure peak_conversion --output /tmp/sweep_torsion

# 2. Extract critical field -- collapses Bpeak dimension
#    Uses Boccaletti formula at B_ref=0.1 to compute threshold
uv run tidal analyze /tmp/sweep_torsion \
  --critical-field Bpeak \
  --reference-formula boccaletti --reference-B 0.1 \
  --output /tmp/critical_torsion

# 3. Plot 1/B_min heatmap (outer params determine plot type)
#    1 outer param -> line plot
#    2 outer params -> 2D heatmap
#    3+ outer params -> parallel coordinates
uv run tidal plot /tmp/critical_torsion --type sweep --metric inv_B_min
\end{lstlisting}

\subsubsection{Theory comparison (amplification factor)}
\label{sec:theory-comparison}

\begin{lstlisting}[style=bash]
# Run both E-M and new theory sweeps with identical parameter grids
uv run tidal sweep examples/data/gertsenshtein_localized.json \
  --sweep "Bpeak=0.01:0.5:30" --sweep "R=1:20:10" \
  --measure peak_conversion --output /tmp/sweep_em

uv run tidal sweep examples/data/torsion_gertsenshtein.json \
  --sweep "Bpeak=0.01:0.5:30" --sweep "R=1:20:10" \
  --measure peak_conversion --output /tmp/sweep_torsion

# Extract B_min for both with the same threshold
uv run tidal analyze /tmp/sweep_em \
  --critical-field Bpeak --threshold 0.99 --output /tmp/critical_em
uv run tidal analyze /tmp/sweep_torsion \
  --critical-field Bpeak --threshold 0.99 --output /tmp/critical_torsion

# Compare the results.csv files:
# amplification(R) = B_min_em(R) / B_min_torsion(R)
\end{lstlisting}

\subsubsection{Simple numeric threshold}
\label{sec:simple-threshold}

\begin{lstlisting}[style=bash]
# Directly specify a threshold (e.g. from literature)
uv run tidal analyze /tmp/sweep_torsion \
  --critical-field Bpeak --threshold 0.5 --output /tmp/critical_simple
\end{lstlisting}

\subsection{CLI Reference}
\label{sec:cf-cli}

\begin{lstlisting}
tidal analyze SWEEP_DIR --critical-field PARAM [options]

Required:
  SWEEP_DIR               Path to sweep output directory
  --critical-field PARAM  Field-strength parameter to threshold on
  --output DIR            Output directory for reduced results

Options:
  --threshold FLOAT       Metric threshold (default: 0.99)
  --metric STR            Metric to threshold on (default: P_final)
  --no-interpolate        Disable linear interpolation
  --reference-formula     {boccaletti,uniform}
                          Compute threshold from E-M analytical formula
  --reference-B FLOAT     Reference B value for analytical formula
\end{lstlisting}

\subsection{Output Files}
\label{sec:cf-output}

The output directory contains standard sweep result files:

\begin{table}[htbp]
\centering
\caption{Output files from critical field analysis.}
\label{tab:cf-output-files}
\begin{tabular}{@{}lp{9cm}@{}}
\toprule
File & Contents \\
\midrule
\texttt{results.csv}  & One row per outer-param combination, with $B_\text{min}$, $1/B_\text{min}$, errors \\
\texttt{results.json} & Same data in JSON format \\
\texttt{sweep.json}   & Provenance metadata (original sweep path, threshold, metric) \\
\bottomrule
\end{tabular}
\end{table}

\subsubsection{Output columns}
\label{sec:cf-columns}

\begin{table}[htbp]
\centering
\caption{Output columns in the critical field results.}
\label{tab:cf-columns}
\begin{tabular}{@{}lp{9cm}@{}}
\toprule
Column & Description \\
\midrule
(outer params)              & Remaining swept parameter values \\
\texttt{B\_min}             & Minimum field strength for conversion $\geq$ threshold \\
\texttt{inv\_B\_min}        & $1/B_\text{min}$ (default heatmap metric --- brighter $=$ easier conversion) \\
\texttt{B\_min\_err}        & Combined error (quadrature of grid + interp + metric) \\
\texttt{B\_min\_err\_grid}  & Error from sweep grid spacing \\
\texttt{B\_min\_err\_metric}& Error from replicate variance (NaN if no replicates) \\
\texttt{B\_min\_err\_interp}& Error from interpolation model \\
\texttt{inv\_B\_min\_err}   & Propagated error on $1/B_\text{min}$ \\
\texttt{crossing\_quality}  & Quality flag: \texttt{good}, \texttt{coarse}, \texttt{edge}, \texttt{none} \\
\bottomrule
\end{tabular}
\end{table}

\subsection{Implementation}
\label{sec:cf-implementation}

\subsubsection{Key files}
\label{sec:cf-key-files}

\begin{table}[htbp]
\centering
\caption{Key implementation files for critical field analysis.}
\label{tab:cf-key-files}
\begin{tabular}{@{}ll@{}}
\toprule
File & Role \\
\midrule
\texttt{tidal/measurement/\_critical\_field.py} & Core algorithm: crossing detection, error estimation, analytical formulas \\
\texttt{tidal/cli/\_analyze.py}                 & CLI dispatch for \texttt{tidal analyze --critical-field} \\
\texttt{tidal/cli/\_\_init\_\_.py}              & CLI argument definitions \\
\texttt{tests/test\_critical\_field.py}          & 32 unit tests covering all edge cases \\
\bottomrule
\end{tabular}
\end{table}

\subsubsection{Architecture}
\label{sec:cf-architecture}

The critical field analysis is a \textbf{post-hoc cross-run reduction}, not a
per-simulation measurement. It operates on completed sweep data:

\begin{lstlisting}
SweepResults (N swept params)
    -> compute_critical_field()
        -> CriticalFieldResult
            -> critical_field_to_sweep_results()
                -> SweepResults (N-1 swept params)
                    -> existing tidal plot infrastructure
\end{lstlisting}

The key design insight: by producing a standard \texttt{SweepResults}, the entire
existing plot infrastructure (1D line, 2D heatmap, parallel coordinates)
works automatically without any plot code changes.

\subsection{Future Work}
\label{sec:cf-future}

\begin{itemize}
  \item \textbf{Arbitrary geometry formulas}: User-provided analytical expressions for
    reference thresholds (tracked in GitHub issue)
  \item \textbf{Paired comparison mode}: \texttt{--reference DIR} to load a reference sweep and
    compute amplification factors directly
  \item \textbf{Transit completeness check}: Detect whether the wavepacket has fully
    exited the B-field region before trusting $P_\text{final}$
  \item \textbf{Replicate-aware $B_\text{min}$}: Compute $B_\text{min}$ per replicate for robust
    uncertainty estimation when ensemble data is available
\end{itemize}

\subsection{References}
\label{sec:cf-references}

\begin{enumerate}
  \item Boccaletti, D.\ et al.\ (1970) Nuovo Cim.\ 70B, 129--146
  \item Dandoy, V.\ et al.~\cite{dandoy2024constraining}
  \item Domcke, V.\ \& Garcia-Cely, C.~\cite{domcke2023simple}
\end{enumerate}
